% Options for packages loaded elsewhere
\PassOptionsToPackage{unicode}{hyperref}
\PassOptionsToPackage{hyphens}{url}
\PassOptionsToPackage{dvipsnames,svgnames,x11names}{xcolor}
%
\documentclass[
  12pt,
]{article}
\usepackage{amsmath,amssymb}
\usepackage{iftex}
\ifPDFTeX
  \usepackage[T1]{fontenc}
  \usepackage[utf8]{inputenc}
  \usepackage{textcomp} % provide euro and other symbols
\else % if luatex or xetex
  \usepackage{unicode-math} % this also loads fontspec
  \defaultfontfeatures{Scale=MatchLowercase}
  \defaultfontfeatures[\rmfamily]{Ligatures=TeX,Scale=1}
\fi
\usepackage{lmodern}
\ifPDFTeX\else
  % xetex/luatex font selection
\fi
% Use upquote if available, for straight quotes in verbatim environments
\IfFileExists{upquote.sty}{\usepackage{upquote}}{}
\IfFileExists{microtype.sty}{% use microtype if available
  \usepackage[]{microtype}
  \UseMicrotypeSet[protrusion]{basicmath} % disable protrusion for tt fonts
}{}
\makeatletter
\@ifundefined{KOMAClassName}{% if non-KOMA class
  \IfFileExists{parskip.sty}{%
    \usepackage{parskip}
  }{% else
    \setlength{\parindent}{0pt}
    \setlength{\parskip}{6pt plus 2pt minus 1pt}}
}{% if KOMA class
  \KOMAoptions{parskip=half}}
\makeatother
\usepackage{xcolor}
\usepackage[margin=1in]{geometry}
\usepackage{color}
\usepackage{fancyvrb}
\newcommand{\VerbBar}{|}
\newcommand{\VERB}{\Verb[commandchars=\\\{\}]}
\DefineVerbatimEnvironment{Highlighting}{Verbatim}{commandchars=\\\{\}}
% Add ',fontsize=\small' for more characters per line
\usepackage{framed}
\definecolor{shadecolor}{RGB}{248,248,248}
\newenvironment{Shaded}{\begin{snugshade}}{\end{snugshade}}
\newcommand{\AlertTok}[1]{\textcolor[rgb]{0.94,0.16,0.16}{#1}}
\newcommand{\AnnotationTok}[1]{\textcolor[rgb]{0.56,0.35,0.01}{\textbf{\textit{#1}}}}
\newcommand{\AttributeTok}[1]{\textcolor[rgb]{0.13,0.29,0.53}{#1}}
\newcommand{\BaseNTok}[1]{\textcolor[rgb]{0.00,0.00,0.81}{#1}}
\newcommand{\BuiltInTok}[1]{#1}
\newcommand{\CharTok}[1]{\textcolor[rgb]{0.31,0.60,0.02}{#1}}
\newcommand{\CommentTok}[1]{\textcolor[rgb]{0.56,0.35,0.01}{\textit{#1}}}
\newcommand{\CommentVarTok}[1]{\textcolor[rgb]{0.56,0.35,0.01}{\textbf{\textit{#1}}}}
\newcommand{\ConstantTok}[1]{\textcolor[rgb]{0.56,0.35,0.01}{#1}}
\newcommand{\ControlFlowTok}[1]{\textcolor[rgb]{0.13,0.29,0.53}{\textbf{#1}}}
\newcommand{\DataTypeTok}[1]{\textcolor[rgb]{0.13,0.29,0.53}{#1}}
\newcommand{\DecValTok}[1]{\textcolor[rgb]{0.00,0.00,0.81}{#1}}
\newcommand{\DocumentationTok}[1]{\textcolor[rgb]{0.56,0.35,0.01}{\textbf{\textit{#1}}}}
\newcommand{\ErrorTok}[1]{\textcolor[rgb]{0.64,0.00,0.00}{\textbf{#1}}}
\newcommand{\ExtensionTok}[1]{#1}
\newcommand{\FloatTok}[1]{\textcolor[rgb]{0.00,0.00,0.81}{#1}}
\newcommand{\FunctionTok}[1]{\textcolor[rgb]{0.13,0.29,0.53}{\textbf{#1}}}
\newcommand{\ImportTok}[1]{#1}
\newcommand{\InformationTok}[1]{\textcolor[rgb]{0.56,0.35,0.01}{\textbf{\textit{#1}}}}
\newcommand{\KeywordTok}[1]{\textcolor[rgb]{0.13,0.29,0.53}{\textbf{#1}}}
\newcommand{\NormalTok}[1]{#1}
\newcommand{\OperatorTok}[1]{\textcolor[rgb]{0.81,0.36,0.00}{\textbf{#1}}}
\newcommand{\OtherTok}[1]{\textcolor[rgb]{0.56,0.35,0.01}{#1}}
\newcommand{\PreprocessorTok}[1]{\textcolor[rgb]{0.56,0.35,0.01}{\textit{#1}}}
\newcommand{\RegionMarkerTok}[1]{#1}
\newcommand{\SpecialCharTok}[1]{\textcolor[rgb]{0.81,0.36,0.00}{\textbf{#1}}}
\newcommand{\SpecialStringTok}[1]{\textcolor[rgb]{0.31,0.60,0.02}{#1}}
\newcommand{\StringTok}[1]{\textcolor[rgb]{0.31,0.60,0.02}{#1}}
\newcommand{\VariableTok}[1]{\textcolor[rgb]{0.00,0.00,0.00}{#1}}
\newcommand{\VerbatimStringTok}[1]{\textcolor[rgb]{0.31,0.60,0.02}{#1}}
\newcommand{\WarningTok}[1]{\textcolor[rgb]{0.56,0.35,0.01}{\textbf{\textit{#1}}}}
\usepackage{graphicx}
\makeatletter
\def\maxwidth{\ifdim\Gin@nat@width>\linewidth\linewidth\else\Gin@nat@width\fi}
\def\maxheight{\ifdim\Gin@nat@height>\textheight\textheight\else\Gin@nat@height\fi}
\makeatother
% Scale images if necessary, so that they will not overflow the page
% margins by default, and it is still possible to overwrite the defaults
% using explicit options in \includegraphics[width, height, ...]{}
\setkeys{Gin}{width=\maxwidth,height=\maxheight,keepaspectratio}
% Set default figure placement to htbp
\makeatletter
\def\fps@figure{htbp}
\makeatother
\setlength{\emergencystretch}{3em} % prevent overfull lines
\providecommand{\tightlist}{%
  \setlength{\itemsep}{0pt}\setlength{\parskip}{0pt}}
\setcounter{secnumdepth}{5}
\ifLuaTeX
  \usepackage{selnolig}  % disable illegal ligatures
\fi
\usepackage{bookmark}
\IfFileExists{xurl.sty}{\usepackage{xurl}}{} % add URL line breaks if available
\urlstyle{same}
\hypersetup{
  pdftitle={Getting Started},
  pdfauthor={Center for Veterinary Biologics - Statistics Section},
  colorlinks=true,
  linkcolor={Maroon},
  filecolor={Maroon},
  citecolor={Blue},
  urlcolor={blue},
  pdfcreator={LaTeX via pandoc}}

\title{Getting Started}
\usepackage{etoolbox}
\makeatletter
\providecommand{\subtitle}[1]{% add subtitle to \maketitle
  \apptocmd{\@title}{\par {\large #1 \par}}{}{}
}
\makeatother
\subtitle{Getting Started}
\author{Center for Veterinary Biologics - Statistics Section}
\date{August 2019}

\begin{document}
\maketitle

{
\hypersetup{linkcolor=}
\setcounter{tocdepth}{3}
\tableofcontents
}
\section{Introduction}\label{introduction}

The \texttt{DiagTestKit} package provides functions to obtain point and
interval estimates for diagnostic sensitivity and specificity for an
experimental diagnostic test kit when one or more imperfect reference
tests are used. The package also includes a function to estimate and
obtain confidence intervals for diagnostic sensitivity or specificity in
the event an infallible reference test is used.

The technical details for how the algorithm works are provided in
\href{https://www.aphis.usda.gov/aphis/ourfocus/animalhealth/veterinary-biologics/biologics-regulations-and-guidance/ct_vb_statwi}{STATWI0002}.
For a 2-state experimental test with one or more fallible reference
tests, the default display will show the point and interval estimates
for sensitivity (\(\pi_1\)) and specificity (\(\theta_1\)). For a
3-state experimental test with one or more fallible reference tests, the
default display will show the point and interval estimates for
sensitivity (\(\pi_1\)), specificity (\(\theta_1\)), probability of
suspect for disease positive (\(\psi_1\)) and probability of suspect for
disease negative (\(\phi_1\)). The algorithm optimizes the sum of
squared residuals with respect to \(\pi_1\), \(\theta_1\), \(\delta_1\),
and \(\gamma_1\) where \(\delta_1\) and \(\gamma_1\) express the
probability of a suspect result as a fraction of the non-correct test
results. See
\href{https://www.aphis.usda.gov/aphis/ourfocus/animalhealth/veterinary-biologics/biologics-regulations-and-guidance/ct_vb_statwi}{STATWI0002}
for further details.

This vignette provides step--by--step instructions for using the package
to obtain point and interval estimates for the sensitivity (Sn =
\(\pi\)) and specificity (Sp = \(\theta\)) of a 2--state experimental
test kit. The \texttt{DiagTestKit} package requires a functional
installation of R software. \texttt{DiagTestKit} package is available
via
\href{https://github.com/ABS-dev/DiagTestKit/blob/master/README.md}{github}.
Code in this vignette was run using version 0.6.11 of the
\texttt{DiagTestKit} package.

\section{Data processing}\label{data-processing}

The data file titled dichotomoussensspec\_deviceinfo.csv will be used to
illustrate the use of the \texttt{estimateSnSp} function. While this
data file includes 3 sample types, only the whole blood samples will be
used to demonstrate the \texttt{estimateSnSp} function. Importing the
data is illustrated for completeness, but the dataset dat\_dichot is
available within the \texttt{DiagTestKit} package. All raw data files
should follow the formats in Section 1.8 or Section 1.9 of Appendix 1 in
\href{https://www.aphis.usda.gov/animal_health/vet_biologics/publications/16-CVBDataGuide.pdf}{CVB
Data Guide}. A table of counts (dat\_infal) available within the
``DiagTestKit'' package will be used to illustrate the use of the
\texttt{cloppearSnSp} function. The \texttt{cloppearSnSp} provides the
binomial confidence interval using the Clopper-Pearson method when an
infallible reference test is used.

\subsection{Load package and data}\label{load-package-and-data}

\begin{Shaded}
\begin{Highlighting}[]
\FunctionTok{library}\NormalTok{(ggplot2)}
\FunctionTok{library}\NormalTok{(DiagTestKit)}
\FunctionTok{data}\NormalTok{(}\StringTok{"ExampleData"}\NormalTok{)}
\NormalTok{dat }\OtherTok{\textless{}{-}}\NormalTok{ dat\_dichot}
\NormalTok{infallible }\OtherTok{\textless{}{-}}\NormalTok{ dat\_infal}
\end{Highlighting}
\end{Shaded}

\subsection{Obtaining counts}\label{obtaining-counts}

The data input for the \texttt{estimateSnSp} function requires creating
a data frame that includes the count for each unique combination of test
results. The data frame should include zero counts for possible test
combinations that were not observed, but it should not include zero
counts for impossible test combinations (structural zeros). Including
non-structural zeros may be achieved using the .drop=FALSE setting
within \texttt{ddply}. There is more than one way to obtain the data
frame with the counts for all possible test combinations. Two
possibilities are illustrated here. First, the \texttt{ddply} function
in the \texttt{plyr} package can be used.

\begin{Shaded}
\begin{Highlighting}[]
\CommentTok{\#Blood Samples}
\NormalTok{blood }\OtherTok{\textless{}{-}} \FunctionTok{ddply}\NormalTok{(}\AttributeTok{.data =} \FunctionTok{subset}\NormalTok{(dat, specimen }\SpecialCharTok{==} \StringTok{"wholeblood"}\NormalTok{),}
                  \AttributeTok{.variables =}\NormalTok{ .(visual\_read, ref\_result),}
                  \AttributeTok{.drop =} \ConstantTok{FALSE}\NormalTok{,}
                  \AttributeTok{.fun =}\NormalTok{ summarize,}
                  \AttributeTok{Count =} \FunctionTok{length}\NormalTok{(deviceID))}
\end{Highlighting}
\end{Shaded}

This will require renaming the columns in a manner that will be
recognized by \texttt{estimateSnSp}. Namely, the column name of
visual\_read from the raw data file needs to be changed to reflect this
is the response associated with the experimental kit. The expression
``exp'' should be included in the column corresponding to the
experimental test results. The expression ``ref'' should be included in
all columns corresponding to the reference test(s) results.

\begin{Shaded}
\begin{Highlighting}[]
\FunctionTok{names}\NormalTok{(blood) }\OtherTok{\textless{}{-}} \FunctionTok{c}\NormalTok{(}\StringTok{"Exp"}\NormalTok{, }\StringTok{"Ref1"}\NormalTok{, }\StringTok{"Count"}\NormalTok{)}
\end{Highlighting}
\end{Shaded}

Alternatively, the data frame including the counts from all possible
testing combinations can be obtained by converting the output from the
\texttt{table} function in \texttt{base} R into a data frame.

\begin{Shaded}
\begin{Highlighting}[]
\NormalTok{blood2 }\OtherTok{\textless{}{-}} \FunctionTok{data.frame}\NormalTok{(}\FunctionTok{table}\NormalTok{(}\AttributeTok{Exp =}\NormalTok{ dat}\SpecialCharTok{$}\NormalTok{visual\_read[dat}\SpecialCharTok{$}\NormalTok{specimen }\SpecialCharTok{==} \StringTok{"wholeblood"}\NormalTok{],}
                          \AttributeTok{Ref1 =}\NormalTok{ dat}\SpecialCharTok{$}\NormalTok{ref\_result[dat}\SpecialCharTok{$}\NormalTok{specimen }\SpecialCharTok{==} \StringTok{"wholeblood"}\NormalTok{]))}
\end{Highlighting}
\end{Shaded}

\section{Fallible Reference Test(s)}\label{fallible-reference-tests}

Use the \texttt{estimateSnSp} function to estimate the sensitivity and
specificity of the 2--state experimental test. The function has 4
required inputs, the data frame of counts (dat), reasonable values based
on available information for the sensitivity of each reference test
(Sn.ref), reasonable values based on available information for the
specificity of each reference test (Sp.ref), and a vector containing
reasonable estimates of prevalence for each population sampled. Each
element in the prevalence vector should be named to distinguish
populations or sample sets (i.e.~a named vector). Two alternative
methods for inputting the performance characteristics (sensitivity and
specificity) of the reference test are available (and illustrated below)
for a 2--stage reference test, a named vector or a named data frame. For
a 3--stage reference test, the probability of suspect as a proportion of
an incorrect test result (for both disease positive and disease
negative) must be supplied. Therefore, if one or more reference tests
have 3 stages, a named data frame must be used.

\subsection{Use a named data frame for inputs of Sn.ref and
Sp.ref}\label{use-a-named-data-frame-for-inputs-of-sn.ref-and-sp.ref}

This example illustrates the use of a named data frame rather than a
named vector for the input variables Sn.ref and Sp.ref. The output seen
here indicates that there was no column named ``population'' and the
dataset is being treated as if it has been sampled from a single
population. The function also provides updates every 50 iterations
(simulation cycles). Here nsim was not specified in the input and
therefore the default of 1000 simulation cycles is being performed.

Here, the dataset blood is used as the input data (dat). Both Sn.ref and
Sp.ref are named data frames containing a single column and two rows.
Because the reference test only has 2 states, the second row includes a
zero (as there is no suspect region). The column is named ``Ref1'' to
indicate the inputs correspond to the first (and only, in this instance)
reference test. The prev.pop is a vector with a single named element.
The element name is ``A'' used as an arbitrary designation for the only
population (or sample set) tested.

\begin{Shaded}
\begin{Highlighting}[]
\NormalTok{blood\_SnSp }\OtherTok{\textless{}{-}} \FunctionTok{estimateSnSp}\NormalTok{(}\AttributeTok{dat =}\NormalTok{ blood,}
                           \AttributeTok{Sn.ref =} \FunctionTok{data.frame}\NormalTok{(}\AttributeTok{Ref1 =} \FunctionTok{c}\NormalTok{(}\FloatTok{0.95}\NormalTok{, }\DecValTok{0}\NormalTok{)),}
                           \AttributeTok{Sp.ref =} \FunctionTok{data.frame}\NormalTok{(}\AttributeTok{Ref1 =} \FunctionTok{c}\NormalTok{(}\FloatTok{0.98}\NormalTok{, }\DecValTok{0}\NormalTok{)),}
                           \AttributeTok{prev.pop =} \FunctionTok{c}\NormalTok{(}\AttributeTok{A =} \FloatTok{0.70}\NormalTok{))}
\end{Highlighting}
\end{Shaded}

\begin{verbatim}
## Warning in estimateSnSp(dat = blood, Sn.ref =
## data.frame(Ref1 = c(0.95, : The data suggests a
## single population was tested
\end{verbatim}

\begin{verbatim}
## The optimization has begun
\end{verbatim}

\begin{verbatim}
## The following is the number of iterations completed: 50
\end{verbatim}

\begin{verbatim}
## The following is the number of iterations completed: 100
\end{verbatim}

\begin{verbatim}
## The following is the number of iterations completed: 150
\end{verbatim}

\begin{verbatim}
## The following is the number of iterations completed: 200
\end{verbatim}

\begin{verbatim}
## The following is the number of iterations completed: 250
\end{verbatim}

\begin{verbatim}
## The following is the number of iterations completed: 300
\end{verbatim}

\begin{verbatim}
## The following is the number of iterations completed: 350
\end{verbatim}

\begin{verbatim}
## The following is the number of iterations completed: 400
\end{verbatim}

\begin{verbatim}
## The following is the number of iterations completed: 450
\end{verbatim}

\begin{verbatim}
## The following is the number of iterations completed: 500
\end{verbatim}

\begin{verbatim}
## The following is the number of iterations completed: 550
\end{verbatim}

\begin{verbatim}
## The following is the number of iterations completed: 600
\end{verbatim}

\begin{verbatim}
## The following is the number of iterations completed: 650
\end{verbatim}

\begin{verbatim}
## The following is the number of iterations completed: 700
\end{verbatim}

\begin{verbatim}
## The following is the number of iterations completed: 750
\end{verbatim}

\begin{verbatim}
## The following is the number of iterations completed: 800
\end{verbatim}

\begin{verbatim}
## The following is the number of iterations completed: 850
\end{verbatim}

\begin{verbatim}
## The following is the number of iterations completed: 900
\end{verbatim}

\begin{verbatim}
## The following is the number of iterations completed: 950
\end{verbatim}

\begin{verbatim}
## The following is the number of iterations completed: 1000
\end{verbatim}

\texttt{\#\#\ The\ optimization\ has\ begun}
\texttt{\#\#\ The\ following\ is\ the\ number\ of\ iterations\ completed:\ 50}
\texttt{\#\#\ The\ following\ is\ the\ number\ of\ iterations\ completed:\ 100}
\texttt{\#\#\ The\ following\ is\ the\ number\ of\ iterations\ completed:\ 150}
\(\vdots\)
\texttt{\#\#\ The\ following\ is\ the\ number\ of\ iterations\ completed:\ 1000}

\begin{Shaded}
\begin{Highlighting}[]
\NormalTok{blood\_SnSp}
\end{Highlighting}
\end{Shaded}

\begin{verbatim}
## 1000  simulations 
## 95 % Interval Estimates
## 
##               Point.Estimate     Lower Upper
## Sn = P(T+|D+)      1.0000000 0.9965680     1
## Sp = P(T-|D-)      0.9505929 0.8403921     1
\end{verbatim}

\subsection{Use a named vector for inputs of Sn.ref and
Sp.ref}\label{use-a-named-vector-for-inputs-of-sn.ref-and-sp.ref}

As this example has a single 2--state reference test, a named vector can
be used for the input variables for Sn.ref and Sp.ref (rather than a
named data frame). If one or more of the reference tests are 3--state
tests (includes a suspect region), using a named vector is not an
option. Using named vectors for Sn.ref and Sp.ref is illustrated with
the data frame obtained by converting the output from the \texttt{table}
function (blood2). This example shows the additional code to reduce the
frequency of messages regarding the number of iterations completed.

\begin{Shaded}
\begin{Highlighting}[]
\NormalTok{blood2\_SnSp }\OtherTok{\textless{}{-}} \FunctionTok{estimateSnSp}\NormalTok{(}\AttributeTok{dat =}\NormalTok{ blood2,}
                           \AttributeTok{Sn.ref =} \FunctionTok{c}\NormalTok{(}\AttributeTok{Ref1 =} \FloatTok{0.95}\NormalTok{),}
                           \AttributeTok{Sp.ref =} \FunctionTok{c}\NormalTok{(}\AttributeTok{Ref1 =} \FloatTok{0.98}\NormalTok{),}
                           \AttributeTok{prev.pop =} \FunctionTok{c}\NormalTok{(}\AttributeTok{A =} \FloatTok{0.70}\NormalTok{),}
                           \AttributeTok{control =} \FunctionTok{estimateSnSpControl}\NormalTok{(}\AttributeTok{iter.n =} \DecValTok{250}\NormalTok{))}
\end{Highlighting}
\end{Shaded}

\begin{verbatim}
## Warning in estimateSnSp(dat = blood2, Sn.ref =
## c(Ref1 = 0.95), Sp.ref = c(Ref1 = 0.98), : The
## data suggests a single population was tested
\end{verbatim}

\begin{verbatim}
## The optimization has begun
\end{verbatim}

\begin{verbatim}
## The following is the number of iterations completed: 250
\end{verbatim}

\begin{verbatim}
## The following is the number of iterations completed: 500
\end{verbatim}

\begin{verbatim}
## The following is the number of iterations completed: 750
\end{verbatim}

\begin{verbatim}
## The following is the number of iterations completed: 1000
\end{verbatim}

\begin{Shaded}
\begin{Highlighting}[]
\NormalTok{blood2\_SnSp}
\end{Highlighting}
\end{Shaded}

\begin{verbatim}
## 1000  simulations 
## 95 % Interval Estimates
## 
##               Point.Estimate     Lower Upper
## Sn = P(T+|D+)      1.0000000 0.9951525     1
## Sp = P(T-|D-)      0.9525117 0.8434030     1
\end{verbatim}

\subsection{Controlling Output}\label{controlling-output}

\subsubsection{Seeds}\label{seeds}

The output from the 2 examples (shown above) is not exactly the same
because the seed for the random number generation was not set in these
examples, but the seed used can be viewed using the following code.

\begin{Shaded}
\begin{Highlighting}[]
\NormalTok{blood\_SnSp}\SpecialCharTok{$}\NormalTok{input}\SpecialCharTok{$}\NormalTok{seed}
\end{Highlighting}
\end{Shaded}

\begin{verbatim}
## [1] 60200
\end{verbatim}

\begin{Shaded}
\begin{Highlighting}[]
\NormalTok{blood2\_SnSp}\SpecialCharTok{$}\NormalTok{input}\SpecialCharTok{$}\NormalTok{seed}
\end{Highlighting}
\end{Shaded}

\begin{verbatim}
## [1] 28937
\end{verbatim}

The output can be forced to be the same by controlling the seed. This
will be illustrated by using the second dataset and the seed used for
the first call to \texttt{estimateSnSp}. The message regarding the
number of iterations has been suppressed in this example.

\begin{Shaded}
\begin{Highlighting}[]
\NormalTok{blood2\_a\_SnSp }\OtherTok{\textless{}{-}}
  \FunctionTok{estimateSnSp}\NormalTok{(}\AttributeTok{dat =}\NormalTok{ blood2,}
      \AttributeTok{Sn.ref =} \FunctionTok{c}\NormalTok{(}\AttributeTok{Ref1 =} \FloatTok{0.95}\NormalTok{),}
      \AttributeTok{Sp.ref =} \FunctionTok{c}\NormalTok{(}\AttributeTok{Ref1 =} \FloatTok{0.98}\NormalTok{),}
      \AttributeTok{prev.pop =} \FunctionTok{c}\NormalTok{(}\AttributeTok{A =} \FloatTok{0.70}\NormalTok{),}
      \AttributeTok{control =} \FunctionTok{estimateSnSpControl}\NormalTok{(}\AttributeTok{seed =}\NormalTok{ blood\_SnSp}\SpecialCharTok{$}\NormalTok{input}\SpecialCharTok{$}\NormalTok{seed,}
                                      \AttributeTok{rep.iter =} \ConstantTok{FALSE}\NormalTok{))}
\end{Highlighting}
\end{Shaded}

\begin{verbatim}
## Warning in estimateSnSp(dat = blood2, Sn.ref =
## c(Ref1 = 0.95), Sp.ref = c(Ref1 = 0.98), : The
## data suggests a single population was tested
\end{verbatim}

\begin{verbatim}
## The optimization has begun
\end{verbatim}

\begin{Shaded}
\begin{Highlighting}[]
\NormalTok{blood2\_a\_SnSp}
\end{Highlighting}
\end{Shaded}

\begin{verbatim}
## 1000  simulations 
## 95 % Interval Estimates
## 
##               Point.Estimate     Lower Upper
## Sn = P(T+|D+)      1.0000000 0.9965680     1
## Sp = P(T-|D-)      0.9505929 0.8403921     1
\end{verbatim}

\subsection{Evaluating Output}\label{evaluating-output}

\subsubsection{Convergence}\label{convergence}

Prior to reporting the point and interval estimates, verify the
optimization convergence in each instance. A numeric value of 0
indicates the optimization converged. Here, we see that all have
converged and we can check the convergence message as well.

\begin{Shaded}
\begin{Highlighting}[]
\FunctionTok{unique}\NormalTok{(blood\_SnSp}\SpecialCharTok{$}\NormalTok{detailOut}\SpecialCharTok{$}\NormalTok{Converge)}
\end{Highlighting}
\end{Shaded}

\begin{verbatim}
## [1] 0
\end{verbatim}

\begin{Shaded}
\begin{Highlighting}[]
\FunctionTok{unique}\NormalTok{(blood\_SnSp}\SpecialCharTok{$}\NormalTok{detailOut}\SpecialCharTok{$}\NormalTok{Message)}
\end{Highlighting}
\end{Shaded}

\begin{verbatim}
## [1] "CONVERGENCE: NORM OF PROJECTED GRADIENT <= PGTOL"
## [2] "CONVERGENCE: REL_REDUCTION_OF_F <= FACTR*EPSMCH"
\end{verbatim}

\subsubsection{Distributions of Optimized
Values}\label{distributions-of-optimized-values}

All optimized values for sensitivity of the experimental kit can be
viewed using a histogram.

\includegraphics[alt={Sn figure}]{C:\Users\TFKent\Documents\CVB_Packages\GitHub - DiagTestKit\vignettes\DiagTestKit-Getting_Started_files/figure-latex/unnamed-chunk-11-1}
\pagebreak

All optimized values for specificity of the experimental kit can be
viewed using a histogram.

\includegraphics[alt={Sp figure}]{C:\Users\TFKent\Documents\CVB_Packages\GitHub - DiagTestKit\vignettes\DiagTestKit-Getting_Started_files/figure-latex/unnamed-chunk-12-1}
\pagebreak

\subsubsection{Simulated Values for Reference
Test}\label{simulated-values-for-reference-test}

Similarly, the simulated values of sensitivity for the reference test
can be extracted using \texttt{blood\_SnSp\$input\$Sn.sims}. The
simulated distribution for the sensitivity of the reference test is
displayed as a histogram.

\includegraphics[alt={Sn reference}]{C:\Users\TFKent\Documents\CVB_Packages\GitHub - DiagTestKit\vignettes\DiagTestKit-Getting_Started_files/figure-latex/unnamed-chunk-13-1}

\pagebreak

The simulated values of specificity for the reference test can be
extracted using \texttt{blood\_SnSp\$input\$Sp.sims}. The simulated
distribution for the specificity of the reference is displayed as a
histogram.

\includegraphics[alt={Sp reference}]{C:\Users\TFKent\Documents\CVB_Packages\GitHub - DiagTestKit\vignettes\DiagTestKit-Getting_Started_files/figure-latex/unnamed-chunk-14-1}

\#Infallible Reference Test

Use the \texttt{cloppearSnSp} function to estimate the sensitivity and
specificity of a 2--state experimental test when an infallible reference
test has been used to determine the true disease status of each sample.
The function has one required input, the data frame of counts (dat).

\subsection{Estimating Sensitivity}\label{estimating-sensitivity}

By default the function will estimate sensitivity.

\begin{Shaded}
\begin{Highlighting}[]
\NormalTok{infal\_Sn }\OtherTok{\textless{}{-}} \FunctionTok{cloppearSnSp}\NormalTok{(}\AttributeTok{dat =}\NormalTok{ infallible)}
\NormalTok{infal\_Sn}
\end{Highlighting}
\end{Shaded}

\begin{verbatim}
## [1] "Sn = P(T+|D+): 0.987013 (95% CI: 0.953876, 0.998423)"
\end{verbatim}

The total number of positive samples tested can be extracted using
\texttt{infal\_Sn\$data\$Total.Positive} and the number of positive
samples that tested positive by the experimental test can be obtained
using \texttt{infal\_Sn\$data\$Test.Positive}. The sensitivity estimate
can be obtained by \texttt{infal\_Sn\$calcVal\$Sn} and the lower and
upper confidence limits can be obtained using
\texttt{infal\_Sn\$calcVal\$Sn.LL} and
\texttt{infal\_Sn\$calcVal\$Sn.UL}, respectively.

\subsection{Estimating Specificity}\label{estimating-specificity}

The function will estimate specificity when \texttt{est.Sn\ =\ FALSE}.

\begin{Shaded}
\begin{Highlighting}[]
\NormalTok{infal\_Sp }\OtherTok{\textless{}{-}} \FunctionTok{cloppearSnSp}\NormalTok{(}\AttributeTok{dat =}\NormalTok{ infallible,}
                         \AttributeTok{est.Sn =} \ConstantTok{FALSE}\NormalTok{)}
\NormalTok{infal\_Sp}
\end{Highlighting}
\end{Shaded}

\begin{verbatim}
## [1] "Sp = P(T-|D-): 0.970297 (95% CI: 0.915643, 0.993832)"
\end{verbatim}

Similarly, the total number of negative samples tested can be extracted
using \texttt{infal\_Sp\$data\$Total.Negative} and the number of
negative samples that tested negative by the experimental test can be
obtained using \texttt{infal\_Sp\$data\$Test.Negative}. The specificity
estimate can be obtained by \texttt{infal\_Sp\$calcVal\$Sp} and the
lower and upper confidence limits can be obtained using
\texttt{infal\_Sp\$calcVal\$Sp.LL} and
\texttt{infal\_Sp\$calcVal\$Sp.UL}, respectively.

\end{document}
